
%\needspace{4\baselineskip}
\color{uniblue}{\subsection{Матрица управляющих воздействий (МУВ)}}
\color{black}

\begin{enumerate}[label=\arabic{section}.\arabic{subsection}.\arabic*, labelsep=4pt, leftmargin=0pt, itemindent=57pt]

\item
Назначение матрицы управляющих воздействий - это объединение по логике «ИЛИ» нескольких входов на одно выходное УВ и его запуска при появлении сигнала на любом из объединённых входов.

\item
Матрица управляющих воздействий представляет собой логику набора сигналов срабатывания функций САЧР, АЧР, ЧАПВ, АОСН, УВ АОСН, АПВН, УВ АПВН, УВ ОН и УВ ЦВН. В устройстве матрица представлена функциональным блоком «МУВ».

\item
В состав функционального блока входит 12 функций формирования выходных воздействий «ЛФ УВ(n)» («ЛФ УВ1» ... «ЛФ УВ12»), где $\{n \in \mathbb{N} \mid 1 \leq n \leq 12\}$. 

\item
Выходные воздействия 1, 2 «жестко» привязаны к УВ СКРМ1 и УВ СКРМ2 соответственно, а с 3 по 12 свободно назначаемые выходные воздействия, назначение которых определяется при конкретном проектировании. На каждое выходное воздействие с помощью уставок SFG2 ... SGF80 может быть назначено до 79 внутренних и внешних сигналов срабатывания функций ПА.

\item
Входные сигналы из логики устройства жестко привязаны к уставкам согласно функциональной схеме (пример: «САЧР: Срабатывание» привязан к уставке SGF2 «Действие на УВ входа 1»). На дискретные входы «Вход 1» ... «Вход 10» возможно назначить сигналы согласно требованиям Заказчика.

\item
Уставкой SGF81 «Режим» контролируется режим работы выходного сигнала:
\begin{itemize}
    \item	в режиме «Импульсный» выходной сигнал будет активен только в течение времени таймера Т1;
    \item	в режиме «Следящий» выходной сигнал будет активен пока не пропадут условия его формирования;
    \item	в режиме «Длительный» выходной сигнал сохраняется в памяти до тех пор, пока не поступит команда сброса от кнопки «Сброс МУВ» или не будет выведена функция из работы.
\end{itemize}

\item
Функциональный блок «МУВ» может быть полностью выведен из работы оперативно путем активации сигнала «ОВ МУВ».

\item
Ввод отдельной функции в работу на этапе параметрирования устройства осуществляется своим программным переключателем SGF1 «Ввод\_функции». Информация о введенной функции отображается выходным сигналом «ЛФ УВ(n): Ввод».

\item
Функция может быть выведена из работы оперативно путем активации сигнала «ОВ ЛФ УВ(n)». Выведенное состояние функции характеризуется активным состоянием сигнала «ЛФ УВ(n): Оперативный вывод».

\item
При активном состоянии входного сигнала «ЛФ УВ(n) на сигн» (ключ «ЛФ УВ(n) на сигн» введен) действие ЛФ УВ(n) переводится на сигнал.

\item
Алгоритм работы функции выполнен в соответствии с рисунком \ref{muv:img1}. В таблице \ref{muv:tbl1} приведены параметры, необходимые для настройки функции.

\vspace{3mm}
\begin{figure}[H]
\centering
\includegraphics[width=1\textwidth,height=1\textheight,keepaspectratio]{\fbpath/209. МУВ/_latex/img/muv.pdf}
\captionof{figure}{Функциональная схема ЛФ УВ}\label{muv:img1}
\end{figure}

\end{enumerate}
